% !TEX root = approx8.tex

\subsection{Filtered and persistence derived Fukaya categories}
\label{s:pdfuk} Let $(X, \omega)$ be a symplectic manifold of one of
the following types:
\begin{enumerate}
\item $X$ is closed.
\item $(X, \omega)$ is symplectically convex at infinity.
\item $(X, \omega = d \lambda)$ is a Liouville manifold with a
  prescribed primitive $\lambda$ of the symplectic structure $\omega$,
  and such that $X$ is symplectically convex at infinity with respect
  to these structure).
\item $(X, \omega = d \lambda)$ is a compact Liouville domain.
\end{enumerate}

In what follows we will work with several variants of Fukaya
categories of $X$, and we will specify below in each case the class of
admissible Lagrangians for this purpose. In each of these cases we
will use the following convention. Unless otherwise stated, whenever
$X$ is a manifold with boundary, the closed Lagrangian submanifolds
$\overline{L} \subset X$ will be implicitly assumed to lie in the {\em
  interior} of $X$.

\subsubsection{The exact case} \label{sb:fuk-ex} Here we assume that
$X$ is Liouville (with a given Liouville form $\lambda$). In order to
obtain a graded theory we add the following assumption: $2c_1(X)=0$,
where $c_1(X)$ stands for the 1'st Chern class of the tangent bundle
of $X$, viewed as a complex vector bundle by endowing $X$ with any
$\omega$-compatible almost complex structure.  We also fix a nowhere
vanishing quadratic complex $n$-form (where
$n = \dim_{\mathbb{C}} X$), namely a nowhere vanishing section
$\Theta$ of the bundle $\Omega^n(X, J)^{\otimes 2}$.

Denote by $\lagex$ the collection of closed, graded, marked, exact
Lagrangian submanifold $L = (\overline{L}, h_L, \theta_L)$. Here
$\overline{L} \subset X$ is a closed $\lambda$-exact Lagrangian
submanifold, which we call the {\em underlying Lagrangian of $L$},
$h_L: \overline{L} \longrightarrow \mathbb{R}$ is a primitive of
$\lambda|_{\overline{L}}$ and
$\theta_L: \overline{L} \longrightarrow \mathbb{R}$ is a grading of
$\overline{L}$, see
e.g.\cite{Se:graded},~\cite[Section~3.2.2.2]{BCZ:tpc} for more
details. Sometimes we will write $\lagex(X)$ or $\lagex_{\lambda}(X)$,
instead of simply $\lagex$, in order to keep track of the additional
data $X$ or $\lambda$.

We can now form the Fukaya category of exact Lagrangians in $X$. We
will follow here a variant of the standard construction
from~\cite{Se:book-fukaya-categ} due to Ambrosioni~\cite{Amb:fil-fuk}
(see also~\cite{BCZ:tpc} for a somewhat simpler implementation in a special case) which
yields a filtered $A_{\infty}$-category. The objects of our Fukaya
category are the elements of $\lagex$. To define the morphisms and
higher structures we need perturbation data. The space of admissible
perturbation data $\mathcal{P}$ is described in detail
in~\cite{Amb:fil-fuk}. Once we fix $p \in \mathcal{P}$ we can define
the Fukaya category $\fuk(\lagex; p)$ as
in~\cite{Amb:fil-fuk}. (Sometimes we will write $\fuk(X, \lagex; p)$
to emphasize the ambient manifold $X$.) For $L_0, L_1 \in \lagex$ we
will sometime denote $\hom_{\fuk(\lagex;p)}(L_0,L_1)$ by
$CF(L_0, L_1; p)$ or  by $\hom(L_0,L_1; p)$. We use here
$\mathbb{Z}_2$ for the coefficients in these $\hom$'s. The grading on
$CF(L_0, L_1; p)$ is defined using the given grading functions
$\theta_{L_0}, \theta_{L_1}$ with which $L_0$, $L_1$ are endowed,
respectively.

\subsubsection{Filtered Fukaya categories in the exact
  case} \label{sbsb:fil-ex} Given two geometrically distinct elements
$L_0, L_1 \in \lagex$ (i.e.~$\overline{L}_0 \neq \overline{L}_1$) and
$p \in \mathcal{P}$, denote by
$H^{L_0,L_1}_p: [0,1] \times X \longrightarrow \mathbb{R}$ the
Hamiltonian function prescribed by $p$ for the Floer datum of
$(L_0, L_1)$. Let $x \in CF(L_0, L_1; p)$ be a generator (i.e.~an
$H^{L_0,L_1}_p$-Hamiltonian chord with endpoints on $L_0, L_1$). Its
action is defined by
\begin{equation} \label{symp-act-1} \mathcal A(x) := \int_0^1
  H_p^{L_0,L_1}(t, x(t))dt - \int_0^1 \lambda(\dot{x}(t)) dt +
  h_{L_1}(x(1)) - h_{L_0}(x(0)).
\end{equation}

We use the action to filter the chain complexes $CF(L_0, L_1; p)$,
$L_0, L_1 \in \lagex$. The grading on $CF(L_0, L_1;p)$ follows the
standard recipe~\cite{Se:graded}, using the grading functions
$\theta_{L_0}, \theta_{L_1}$.

In case $\overline{L}_0 = \overline{L}_1$ the Floer complex
$CF(L_0, L_1; p)$ coincides (up to a shift in grading) with the Morse
complex of $\overline{L}:= \overline{L}_0 = \overline{L}_1$ and we set
the action level $\mathcal{A}(x)$ of all the (non-zero) elements
$x \in CF(L_0, L_1; p)$ to be $h_{L_1} - h_{L_0}$ (note that the
latter is a constant since $\overline{L}_0 = \overline{L}_1$). The
(cohomological) grading of a generator $x \in CF(L_0, L_1; p)$ is
defined as $|x| = (n- \textnormal{ind}_f(x)) + \kappa$, where
$f: \overline{L} \longrightarrow \mathbb{R}$ is the Morse function on
$\overline{L}$ presribed by $p$, $\textnormal{ind}_f(x)$ is the Morse
index of the critical point $x$, $n := \dim \overline{L}$, and
$\kappa = \theta_{L_1} - \theta_{L_0}$ (which is again an integer).

The work of~\cite{Amb:fil-fuk} introduces a special class
$\mathcal{P}$ of perturbation data such that for every
$p \in \mathcal{P}$, the Fukaya category $\fuk(\lagex;p)$ is a genuine
filtered and strictly unital $A_{\infty}$-category.  We endow the
category $\fuk(\lagex;p)$ with translation and shift functors as
follows. Let $L = (\overline{L}, h_L, \theta_L) \in \lagex$. We define
the translation and shift functors, respectively, on objects by:
\begin{equation} \label{eq:T-Sigma-ex-1} TL := (\overline{L}, h_L,
  \theta_L+1), \quad \Sigma^r L := (\overline{L}, h_L+r, \theta_L), \;
  \forall r \in \mathbb{R}.
\end{equation}
These extend to maps on $\hom_{\fuk(\lagex;p)}$'s in the obvious way
and furthermore to filtered $A_{\infty}$-functors
$\fuk(\lagex;p) \longrightarrow \fuk(\lagex;p)$ with trivial higher
order terms.  With these definitions we have natural isomorphisms:
\begin{equation} \label{eq:T-Sigma-ex-2}
  \begin{aligned}
    & CF(T^kL_0, T^lL_1; p)^i \cong
    CF(L_0, L_1; p)^{i-k+l}, \; \forall \; k,l \in \mathbb{Z}, \\
    & CF^{\alpha}(\Sigma^rL_0, \Sigma^sL_1;p) \cong
    CF^{\alpha+r-s}(L_0, L_1; p), \; \forall \; \alpha, r, s \in
    \mathbb{R}.
  \end{aligned}
\end{equation}
The superscripts $i, i-k+l$ in the first line stand for degrees and
the superscripts $\alpha, \alpha+r-s$ in the second line are action
levels.  In the first line we have used cohomological grading
conventions. In homological grading conventions the isomorphism from
the first line takes the following form:
$CF(T^kL_0, T^lL_1; p)_j \cong CF(L_0, L_1; p)_{j+k-l}, \; \forall \;
k,l \in \mathbb{Z}$. In case no confusion between grading and  action levels may arise
we will sometimes also write $CF^i(L_0, L_1;p)$ instead of
$CF(L_0, L_1; p)^i$ and similarly in homology.

Before we proceed we remark that the filtered Fukaya category
$\fuk(\lagex;p)$ depends on the choice of $\lambda$ in two
ways. Firstly, the set of objects $\lagex$ depends on $\lambda$ and
secondly the filtrations on the hom's depend  on $\lambda$.
Whenever we need to emphasize the dependence on $\lambda$ we will
write $\lagex_{\lambda}$ for the set of objects and denote this
category by $\fuk(\lagex_{\lambda}, \lambda ;p)$. (The second
$\lambda$ in the notation indicates that the action is measured with
respect to $\lambda$.) Note that if $\lambda' = \lambda + dG$, where
$G: X \longrightarrow \mathbb{R}$ is a smooth function, compactly
supported in the interior of $X$, then there is a canonical
isomorphism of filtered $A_{\infty}$-categories
\begin{equation} \label{eq:can-iso-fuk-lam} \fuk(\lagex_{\lambda},
  \lambda;p) \longrightarrow \fuk(\lagex_{\lambda'}, \lambda';p).
\end{equation}
Its action on objects is given by
$(\overline{L}, h, \theta) \longmapsto (\overline{L},
h+G|_{\overline{L}}, \theta)$. Its action on morphisms is induced by
the identity map and it has zero higher order terms. We will often
identify all the categories $ \fuk(\lagex_{\lambda'}, \lambda';p)$
with $\lambda'$ as above.

However if $\lambda'$ differs form $\lambda$ by a non-exact closed
form then the collection $\lagex$ entirely changes, so in that case
$\fuk(\lagex_{\lambda}, \lambda ;p)$ and
$\fuk(\lagex_{\lambda'}, \lambda' ;p)$ have completely different
objects and there is no general way to compare these categories.

We will discuss further important properties of the categories
$\fuk(\lagex;p)$ in~\S\ref{sb:sys-fuk} below.

\subsubsection{The monotone case} \label{sb:fuk-mon} In case $X$ is
closed, or symplectically convex at infinity we will consider monotone
Lagrangian submanifolds. (Note that the existence of monotone
Lagrangians in $X$ is not obvious, there are  global
obstructions, e.g.~that the ambient manifold $X$ is itself monotone.)

To obtain a graded Floer theory we follow Seidel's approach for graded
Lagrangians submanifolds~\cite{Se:graded}. Fix
$N \in \mathbb{Z}_{\geq 2}$ for which the image of $2c_1(X)$ in
$H^2(X; \mathbb{Z}_N)$ vanishes (we assume that such an $N$ does
exist). Fix an $N$-fold Maslov covering
$Gr^N_{\text{Lag}}(X, \omega) \longrightarrow Gr_{\text{Lag}}(X,
\omega)$ of the Lagrangian Grassmannian bundle of $(T(X), \omega)$.

Denote by $\Lambda$ the Novikov field with coefficients in
$\mathbb{Z}_2$, namely:
\begin{equation} \label{eq:Nov-ring} \Lambda = \Bigl\{
  \sum_{k=0}^{\infty} a_k T^{\lambda_k} \mid a_k \in \mathbb{Z}_2,
  \lim_{k \to \infty} \lambda_k = \infty \Bigr\},
\end{equation}
and let $\Lambda_0 \subset \Lambda$ be the positive Novikov ring:
\begin{equation} \label{eq:Nov-ring-0} \Lambda_0 = \Bigl\{
  \sum_{k=0}^{\infty} a_k T^{\lambda_k} \mid a_k \in \mathbb{Z}_2,
  \lambda_k\geq 0, \lim_{k \to \infty} \lambda_k = \infty \Bigr\}.
\end{equation}

Denote by $\lagmon$ the collection of closed, graded, monotone
Lagrangian submanifolds whose Maslov-$2$ pseudo-holomorphic disk count (through a generic point and for a generic $J$ that tames $\omega$; see for instance \cite{Bi-Co-Sh:LagrSh} \S 3.5 page 66 for a definition) is
$\mathbf{d} \in \Lambda_0$. We write these elements as triples
$L = (\overline{L}, a_L, s_L)$, where:
\begin{itemize}
\item $\overline{L} \subset X$ is a monotone Lagrangian submanifold
  with Maslov-$2$ disk count $\mathbf{d}$.
\item $s_L$ is a $Gr^N_{\text{Lag}}(X, \omega)$-grading of $L$. More
  specifically,
  $s_L: \overline{L} \longrightarrow Gr^N_{\text{Lag}}(X, \omega)$ is
  a lift of the canonical section
  $\overline{L} \longrightarrow Gr_{\text{Lag}}(X,
  \omega)$. See~\cite{Se:graded} for more details.
\item $a_L \in \mathbb{R}$ is a real number that will be used for
  action filtration purposes in~\S\ref{sbsb:fil-mon} below.
\end{itemize}


\subsubsection{Filtered Fukaya categories in the monotone
  case} \label{sbsb:fil-mon} Let
$L_0 = (\overline{L}_0, a_{L_0}, s_{L_0})$,
$L_1 = (\overline{L}_1, a_{L_1}, s_{L_1})$ be two geometrically
distinct elements in $\lagmon$. Let $p \in \mathcal{P}$, and
$H^{L_0,L_1}_p: [0,1] \times X \longrightarrow \mathbb{R}$ the
Hamiltonian function prescribed by $p$ for the Floer datum of
$(L_0, L_1)$. Recall that the Floer complex $CF(L_0, L_1;p)$ is a free
$\Lambda$-module generated by the $H^{L_0,L_1}_p$-Hamiltonian chords
$x$ with endpoints on $\overline{L}_0, \overline{L}_1$.

Let $0 \neq P(T) \in \Lambda$ and $x$ be an
$H^{L_0,L_1}_p$-Hamiltonian chord as above. We define the action of
$P(T)x$ to be:
\begin{equation} \label{eq:action-mon} \mathcal{A}\bigl(P(T)x\bigr) :=
  -\lambda_0 + \int_0^1 H_p^{L_0,L_1}(t, x(t))dt + a_{L_1} - a_{L_0},
\end{equation}
where $\lambda_0 \in \mathbb{R}$ is the minimal exponent that appears
in the formal power series of $P(T) \in \Lambda$, i.e.
$P(T) = a_0 T^{\lambda_0} + \sum_{i=1}^{\infty} a_i T^{\lambda_i}$
with $a_0 \neq 0$ and $\lambda_i > \lambda_0$ for every $i\geq 1$.  We
now extend $\mathcal{A}$ to $CF(L_0,L_1;p)$ as follows. For a
non-trivial element
$c = P_1(T)x_1 + \cdots + P_l(T)x_l \in CF(L_0, L_1;p)$ we define
$$\mathcal{A}(c) = \max \{ \mathcal{A}(P_k(T) x_k) \mid 1\leq k \leq
l\},$$ and finally we set $\mathcal{A}(0) = -\infty$. The filtration
on $CF(L_0,L_1;p)$ is then defined by:
$$CF^{\alpha}(L_0,L_1;p) := 
\{c \in CF(L_0,L_1; p) \mid \mathcal{A}(c) < \alpha \}.$$ Standard
arguments show that the Floer differential $\mu_1$ preserves this
filtration. It is important to note however, that the filtration
levels $CF^{\alpha}(L_0,L_1;p)$ of $CF(L_0,L_1;p)$ are {\em not}
$\Lambda$-modules but rather $\Lambda_0$-modules.

The grading on $CF(L_0, L_1;p)$, when
$\overline{L}_0 \neq \overline{L}_1$, is defined using the $N$-fold
maslov covering $Gr^N_{\text{Lag}}(X, \omega)$ and the gradings
$s_{L_0}, s_{L_1}$ following the recipe from~\cite{Se:graded}. In
contrast to the case described in~\S\ref{sbsb:fil-ex}, here we only
obtain a $\mathbb{Z}_N$-grading.

As in the exact case, following~\cite{Amb:fil-fuk} there is a class of
perturbation data $\mathcal{P}$ such that for every
$p \in \mathcal{P}$ the Fukaya category $\fuk(\lagmon;p)$ is a genuine
filtered and strictly unital $A_{\infty}$-category. The total category
is $\Lambda$-linear, however when viewed as a filtered
$A_{\infty}$-category, $\fuk(\lagmon;p)$ is only $\Lambda_0$-linear.

Similarly to the exact case, here too we have translation and shift
functors. For the translation functor, recall from~\cite{Se:graded} that the $N$-fold
Maslov covering
$Gr^N_{\text{Lag}}(X, \omega) \longrightarrow Gr_{\text{Lag}}(X,
\omega)$ comes with a $\mathbb{Z}_N$-action
$\rho: \mathbb{Z}_N \longrightarrow
\text{Aut}\Bigl(Gr^N_{\text{Lag}}(X, \omega) \longrightarrow
Gr_{\text{Lag}}(X, \omega) \Bigr)$.  We now
define for $L = (\overline{L}, a_L, s_L) \in \lagmon$:
\begin{equation} \label{eq:T-Sigma-mon-1} TL := (\overline{L}, a_L,
  \rho(1) \circ s_L), \quad \Sigma^r L := (\overline{L}, a_L+r, s_L),
  \; \forall r \in \mathbb{R}.
\end{equation}
Both $T$ and $\Sigma^r$ extend to filtered $A_{\infty}$-functors
$\fuk(\lagmon;p) \longrightarrow \fuk(\lagmon;p)$ and the identities
from~\eqref{eq:T-Sigma-ex-2} continue to hold (with the grading being
taken modulo $N$).


\subsubsection{The weakly exact case} \label{sb:fuk-wex} Assume that
$X$ is closed, or symplectically convex at infinity. we will consider
in this case also weakly exact closed Lagrangian submanifolds
$\overline{L} \subset X$. By weakly exact we mean that
$\langle [\omega], A \rangle = 0$ for all classes
$A \in \textnormal{image\,} (\pi_2(X,L) \longrightarrow
H_2(X,L))$.  A necessary condition for such Lagrangians
to exist is that $X$ is symplectically aspherical,
i.e.~$\langle [\omega], A \rangle = 0$ for all spherical classes
$A \in \textnormal{image\,} (\pi_2(X) \longrightarrow H_2(X))$. Also
note that if $(X, \omega)$ is exact then every exact Lagrangian is
automatically weakly exact, regardless of the chosen primitive
$\lambda$ of $\omega$.)

Denote by $\lagwex$ the collection of triples
$L = (\overline{L}, a_L, \theta_L)$, where $\overline{L} \subset X$ is
a closed weakly exact Lagrangian submanifold, $a_L \in \mathbb{R}$ and
$\theta_L: \overline{L} \longrightarrow \mathbb{R}$ is a grading of
$\overline{L}$ as in~\S\ref{sb:fuk-ex}. Similarly to the monotone
case, we will work here with coefficients in the Novikov ring
$\Lambda$ and its positive version $\Lambda_0$. The Floer complexes
$CF(L_0, L_1; p)$ are defined exactly as in the monotone
case~\S\ref{sbsb:fil-mon}. For the action filtration we follow the
recipe from the monotone case, while for the grading we use the recipe
described in the exact case. More specifically, the action filtration
is defined by~\eqref{eq:action-mon} and the grading is defined as
in~\S\ref{sbsb:fil-ex}. The defintion of the translation functor is
the same as in~\eqref{eq:T-Sigma-ex-1} and the shift functors
$\Sigma^r$ are defined as in~\eqref{eq:T-Sigma-mon-1}. The identities
from~\eqref{eq:T-Sigma-ex-2} continue to hold.



\subsubsection{Systems of Fukaya categories and their
  TPC's} \label{sb:sys-fuk} Let $(X, \omega)$ be a symplectic manifold
as at the beginning of~\S\ref{s:pdfuk} and let $\lag$ be a subset of
one of the collections $\lagex$, $\lagmon$ or $\lagwex$ (depending on
the type of $X$). We assume that $\lag$ is closed under all
translations and all shifts.

Let $\mathcal{P}$ be the space of perturbation data, as
in~\S\ref{sbsb:fil-ex} and~\S\ref{sbsb:fil-mon}, and denote for every
$p \in \mathcal{P}$ by $\fuk(\lag;p)$ the full
$A_{\infty}$-subcategory (of the respective Fukaya categories
$\fuk(\lagex;p)$, $\fuk(\lagmon;p)$ or $\fuk(\lagwex;p)$) with the
objects $\lag$. The categories $\fuk(\lag;p)$ inherit the structure of
filtered $A_{\infty}$-categories from their ambient categories.

We will further restrict the space of perturbation data $\mathcal{P}$
to satisfy the following condition. Fix a constant $B > 0$, and
consider only perturbation data $p \in \mathcal{P}$ such that for
every pair of geometrically distinct objects $L_0, L_1 \in \lag$ with
(i.e.~$\overline{L}_0 \neq \overline{L}_1$) we have
\begin{equation} \label{eq:H-smaller-B} H_p^{L_0,L_1}(t,x) \leq B, \;
  \forall \; t \in [0,1], x \in X.
\end{equation}
Recall from \cite{Amb:fil-fuk} that for every $p \in \mathcal{P}$ and $L_0, L_1$ as above, the
function $H_p^{L_0, L_1}$ is strictly positive. As we will see later,
the particular choice of the constant $B$ will not play an important
role for our considerations (intuitively, it should be a small
number), and its main purpose is to have a uniform bound on all the
Hamiltonian functions that appear in the perturbation data. By a
slight abuse of notation we will denote the space of perturbation data
satisfying the additional condition~\eqref{eq:H-smaller-B} also by
$\mathcal{P}$.

Define now $\nu: \mathcal{P} \longrightarrow \mathbb{R}_{>0}$ by:
\begin{equation} \label{eq:nu-fuk}
  \nu(p) := \sup \Bigl\{H_p^{L_0, L_1}(t,x) \bigm| L_0, L_1 \in \lag
  \text{\ are geometrically distinct, }
  (t,x) \in [0,1]\times X\Bigr\}.
\end{equation}

Next, define a preorder on $\mathcal{P}$. For $p, q \in \mathcal{P}$
we define $p \preceq q$ if for every two geometrically distinct
$L_0, L_1 \in \lag$ we have
$$H_q^{L_0, L_1}(t,x) \leq H_p^{L_0, L_1}(t,x), \; \forall \; (t,x) \in
[0,1] \times X.$$ Note that when endowed with $\preceq$, the set
$\mathcal{P}$ becomes directed (i.e.~in addition to being a preorder
we have that
$\forall \; p', p'' \in \mathcal{P}, \; \exists \; q \in \mathcal{P}$
with $p', p'' \preceq q$).

Let $p, q \in \mathcal{P}$ with $p \preceq q$. From ~\cite{Amb:fil-fuk},
there exists a distinguished class of filtered, strictly unital,
continuation $A_{\infty}$-functors
$\mathcal{H}_{p,q}: \fuk(\lag;p) \longrightarrow
\fuk(\lag;q)$. Different functors in this class are $0$-homotopic
Moreover, we also have a distinguished class of strictly unital
continuation $A_{\infty}$-functors
$\mathcal{H}_{q,p}: \fuk(\lag;q) \longrightarrow \fuk(\lag;p)$ with
linear deviation rate $\leq 2(\nu(p) - \nu(q))$ (see~\S\ref{a:func-LD}
for the precise definition). Here too, any two functors in this class
are $0$-homotopic (the homotopy being a homotopy of LD-functors of a
given deviation rate; see~\S\ref{a:homotopy}).

By~\cite{Amb:fil-fuk} the functors $\mathcal{H}_{p,q}$ staisfy the
conditions listed in~\S\ref{sb:sys-ainfty}, hence the family of
categories $\{\fuk(\lag;p)\}_{p \in \mathcal{P}}$ together with the
comparison functors $\mathcal{H}_{p,q}$, $\mathcal{H}_{q,p}$,
$p \preceq q$, become a system of filtered $A_{\infty}$-categories
which we denote by $\widehat{\fuk}(\lag)$.

The system $\widehat{\fuk}(\lag)$ has in fact two additional important
properties. The first is that all the categories $\fuk(\lag;p)$ have
the same set of objects $\lag$ and the comparison functors act like
the identity on this set. In the terminology of~\S\ref{sbsb:sys-base},
the system $\widehat{\fuk}(\lag)$ has a fixed full base of objects.
The second property is that $\widehat{\fuk}(\lag)$ is in fact a
homotopy system, according to the definitions
from~\S\ref{sbsb:homotopy-sys}. This follows immediately from the
discussion above on the homotopy properties of the comparision
functors.

Note however that although $\widehat{\fuk}(\lag)$ has a fixed full
base of objects, this does not hold anymore for the associated system
of filtered modules $F\md(\widehat{\fuk}(\lag))$ or the Yoneda system
$\mathcal{Y}(\widehat{\fuk}(\lag))$ (in particular, they cannot even
be considered as homotopy systems). This has nothing to do with Fukaya
categories and the reason is purely algebraic - see
Remark~\ref{r:coh-base}. Nevertheless, the fact that
$\widehat{\fuk}(\lag)$ is a homotopy system is still useful, as it
allows to define some of the limit invariants introduced
in~\S\ref{sb:sys-inv}, e.g.~the persistence Hochschild homology
$PHH(\widehat{\fuk}(\lag))$ of the system $\widehat{\fuk}(\lag)$ as
defined in~\eqref{eq:HH-colim}.

Returning to $\widehat{\fuk}(\lag)$, we will turn it into a genuine
system of TPC's with increasing accuracy. There are two ways to do
this, each leading to a slightly different system. The first way is to
use the persistence derived Fukaya category using filtered
modules as described in~\S\ref{sbsb:fmod}, and the second one is to
use the construction from~\S\ref{sbsb:fmod-2}.

We begin with the first implementation. We use the parameter space
$\mathcal{P}$, the size function
$\nu: \mathcal{P} \longrightarrow \mathbb{R}_{>0}$ and the pre-order
$\preceq$ discussed above. The categories in the system of TPC's will
be the persistence derived categories $PD(\fuk(\lag; p))$,
$p \in \mathcal{P}$. The comparison functors
for $p \preceq q$ are defined to be
$$\msh_{p,q} := PD (\mathcal{H}_{p,q}): PD(\fuk(\lag; p))
\longrightarrow PD(\fuk(\lag; q)),$$ where we use here the notation
and construction from~\S\ref{sbsb:fmod}. For $p \preceq q$ we define
$\msh_{q,p}: PD(\fuk(\lag; q)) \longrightarrow PD(\fuk(\lag;
p))$ in the same way. Note that although $\mathcal{H}_{q,p}$ has
linear deviation, its induced functor on the persistence derived
categories is still a TPC-functor. See~\S\ref{sbsb:fmod} for more
details. We will denote this system of TPC's by
$PD(\widehat{\fuk}(\lag))$.

The second implementation is very similar, only that it appeals to the
construction from~\S\ref{sbsb:fmod-2}. The categories in the system
are now $PD^c(\fuk(\lag;p))$, $p \in \mathcal{P}$, and the comparison
functors are $\msh^c_{p,q} := PD^c(\mathcal{H}_{p,q})$ and similarly
for $\msh^c_{q,p}$. We denote this system of TPC's by
$PD^c(\widehat{\fuk}(\lag))$.

Finally, we remark that although $\widehat{\fuk}(\lag)$ is a system
with a fixed full base of objects (and even a homotopy system),
neither $PD(\widehat{\fuk}(\lag))$ nor $PD^c(\widehat{\fuk}(\lag))$
has a coherent full base of objects.

\subsubsection{Functors associated with
  symplectomorphisms} \label{sbsb:symplecto-fun} Let $(X, \omega)$ be
a symplectic manifold of one of the types listed at the beginning
of~\S\ref{s:pdfuk} and $\lag$ be a subset of the collection of
Lagrangians in $X$ as at the beginning of~\S\ref{sb:sys-fuk}. Denote
by $\mathcal{P}$ the space of admissible perturbation data
parametrizing the system of filtered Fukaya categories with objects
$\lag$.

Let $\phi : X \longrightarrow X$ be a symplectic diffeomorphism
compactly supported in the interior of $X$. We will define below an endo-functor of each of the systems of TPC's $PD(\widehat{\fuk}(\lag))$ and
$PD^c(\widehat{\fuk}(\lag))$, that is induced by $\phi$.

In order to define such a functor we need several additional
assumptions on $\phi$. First of all we assume that $\phi$ preserves
the collection of underlying Lagrangians $\overline{L}$ with
$L \in \lag$. The other assumptions on $\phi$ depend on the collection
of Lagrangians (exact, monotone or weakly exact) which $\lag$ is
associated with.

We begin with the case of exact Lagrangians. Here we first need to
keep in mind the Liouville form $\lambda$. Let
$\lag_{\lambda} \subset \lagex_{\lambda}$ be a subset which is closed
under translations and shifts. We will make the following additional
assumptions: $\phi$ is assumed to be exact, namely
$\phi^*\lambda = \lambda + dF$ for some smooth function
$F: X \longrightarrow \mathbb{R}$ which is compactly supported inside
the interior of $X$. We also assume that $\phi$ preserves the homotopy
class of the quadratic complex $n$-form $\Theta$ (among non-vanishing
forms of this type, and after indentifying $(T(X),J)$ with
$(T(X), \phi^*J)$). We also assume that $\phi$ is graded
(see~\cite{Se:graded} for the definition of a graded
symplectomorphism) and denote its action on graded Lagrangians by
$(\overline{L}, \theta) \longmapsto (\phi(\overline{L}),
\phi_*(\theta))$.

Under the above assumptions we obtain a bijective map,
\begin{equation} \label{eq:phi-obj-lam-1} \Ob(\fuk(\lag, \lambda; p))
  \longrightarrow \Ob(\fuk(\lag, \lambda';q)), \quad (\overline{L}, h,
  \theta) \longmapsto (\phi(\overline{L}), h \circ
  \phi^{-1}|_{\phi(\overline{L})}, \phi_*(\theta)),
\end{equation}
where $\lambda' := (\phi^{-1})^*\lambda$. Here $p,q$ are any two
choices of perturbation data (indeed, below we will need to work with
$p \neq q$).

By the assumptions on $\phi$ we have
$\lambda' = \lambda - d(F \circ \phi^{-1})$, hence after composing the
preceding map with the canonical
isomorphism~\eqref{eq:can-iso-fuk-lam} we obtain a bijection
\begin{equation} \label{eq:phi-obj}
  \begin{aligned}
    & \phi: \Ob(\fuk(\lag, \lambda;p)) \longrightarrow \Ob(\fuk(\lag,
    \lambda;q)),
    \\
    & (\overline{L}, h, \theta) \longmapsto (\phi(\overline{L}), h
    \circ \phi^{-1}|_{\phi(\overline{L})} - F|_{\overline{L}} \circ
    \phi^{-1}|_{\phi(\overline{L})}, \phi_*(\theta)),
  \end{aligned}
\end{equation}
which by abuse of notation we still denote by $\phi$.

Next we define a map
$\phi_* : \mathcal{P} \longrightarrow \mathcal{P}$ by pushing forward
the structures in the perturbation datum. Specifically, let
$p \in \mathcal{P}$ be a perturbation data. Recall that $p$ assigns to
each tuple of Lagrangians in $\lag$ some auxiliary structures.  In the
simplest case, when this tuple $\vec{L} = (L_0, \ldots, L_d)$ has no
consecutive entries (in the cyclic sense) with the same underlying
Lagrangians, the perturbation $p$ assigns to $\vec{L}$ a tuple
$(K^p, J^p)$ with two components: the first is a family
$K^p =\{K^p_S \in \Omega^1(S; C_0^{\infty}(X))\}_{S \in \mathcal{S}}$
of $1$-forms parametrized by the space $\mathcal{S}$ of
boundary-punctured disks (with any number $\geq 2$ of punctures). Each
member of the family $K^p$ is a $1$-form $K_S^p$ on the punctured disk
$S$ with values in the space of compactly supported smooth functions
on $X$ (which should be viewed as Hamiltonian functions). The second
component $J^p$ is a family $\{J_S^p\}_{S \in \mathcal{S}}$ of (domain
dependent) almost complex structure, or more precisely each $J^p_S$ is
itelsf an $S$-parametrized family $\{J^p_{S,z}\}_{z \in S}$ of
$\omega$-compatible almost complex structures on $X$ which are also
compatible with the symplectic convexity of $(X, \omega)$ at infinity
(in case $X$ is not closed).  The push forward $\phi_*(p)$ of $p$ by
$\phi$ assigns to $\phi(\vec{L}) := (\phi(L_0), \ldots, \phi(L_d))$
the perturbation data $\phi_*(p) = (\phi_* K^p, \phi_* J^p)$, where
\begin{equation} \label{eq:push-forward}
%  \begin{aligned}
  (\phi_*K^p)_S (\xi) :=  K^p_S(\xi) \circ \phi^{-1}, \; \forall
  \xi \in T(S), \quad (\phi_* J^p)_{S,z} := D \phi \circ J_{S,z}
  \circ D \phi^{-1}, \; \forall S \in \mathcal{S}, z \in S.
%  \end{aligned}
\end{equation}
This completes the definition of $\phi_*(p)$ in the case when
$\vec{L}$ does not contain consecutive entries whose underlying
Lagrangians coincide.

When the tuple of Lagrangians does contain consecutive entries with
the same underlying Lagrangians the perturbation data $p$ are defined
over the space clusters (see~\cite{Amb:fil-fuk}, and
also~\cite[Section~3.3]{BCZ:tpc} where these are called decorated
clusters of punctured disks). In essence these combine punctured disks
and some trees attached to some of the boundary arcs of the punctured
disks. Over each punctured disk in the cluster the perturbation data
are as before (namely, they consits of a pair $(K^p,J^p)$ as described
above). Over each tree the perturbation data consists of a Morse data,
which is a collection of paramter-depending Morse functions and
Riemannian metrics assigned to each edge of the graph (and depending
on a parameter varying along each edge). The push forward by $\phi$
of the Morse data (in the tree part of a cluster) is simply done by
composing the Morse functions with $\phi^{-1}$ and similarly for the
Riemannian metrics. Finally, over each punctured disk in a cluster we
use the same recipe as in~\eqref{eq:push-forward}.

The space of admissible perturbation data
$\mathcal{P}$ is closed under the action of $\phi$, namely for every
$p \in \mathcal{P}$ we have $\phi_*(p) \in \mathcal{P}$. This follows
from the definition of $\mathcal{P}$ in~\cite{Amb:fil-fuk}. Moreover,
it follows from~\eqref{eq:nu-fuk} that for every $p \in \mathcal{P}$ we
have $\nu(\phi_*(p)) = \nu(p)$ and that for every $p \preceq q$ we
have $\phi_*(p) \preceq \phi_*(q)$.

The diffeomorphism $\phi$ also induces obvious chain isomorphisms,
defined for every $L_0, L_1 \in \lag_{\lambda}$ and
$p \in \mathcal{P}$:
\begin{equation} \label{eq:phi-chain} \phi^{CF} : CF(L_0, L_1;
  \lambda; p) \longrightarrow CF(\phi(L_0), \phi(L_1); \lambda;
  \phi_*(p)).
\end{equation}
Note that we are using here the same Liouville form $\lambda$ both in
the domain and target of $\phi^{CF}$, as we have done also
in~\eqref{eq:phi-obj}. Since
$(\phi^{-1})^*\lambda = \lambda - d(F \circ \phi^{-1})$ it follows
that the chain map $\phi^{CF}$ preserves action filtrations.

The maps $\phi^{CF}$ extend to a filtered $A_{\infty}$-functor (in
fact an $A_{\infty}$-isomorphism)
\begin{equation} \label{eq:phi-fun-1} \phi^{\fuk}_p: \fuk(\lag,
  \lambda; p) \longrightarrow \fuk(\lag, \lambda; \phi_*(p))
\end{equation}
whose first order terms are $\phi^{CF}$ and with vanishing higher
order terms. These functors fit together to a functor of systems of
filtered $A_{\infty}$-categories,
$$\widehat{\phi}^{\fuk}: \widehat{\fuk}(\lag, \lambda) \longrightarrow
\widehat{\fuk}(\lag, \lambda),$$ as defined at the end
of~\S\ref{sb:sys-ainfty} (see also~\S\ref{sbsb:func-sys-pc}).

Passing to the persistence derived level, we obtain by~\S\ref{sb:pdc}
TPC functors:
\begin{equation} \label{eq:phi-fun-2} PD(\phi^{\fuk}_p): PD(\fuk(\lag,
  \lambda, p)) \longrightarrow PD(\fuk(\lag, \lambda, \phi_*(p))).
\end{equation}

The collection of functors $PD(\phi_p^{\fuk})$, $p \in \mathcal{P}$,
and the push forward map
$\phi_*: \mathcal{P} \longrightarrow \mathcal{P}$, form together one
functor of systems of TPC's, according to the definitions
from~\S\ref{sb:sys-tpc}:
\begin{equation} \label{eq:phi-fun-3} PD(\widehat{\phi}^{\fuk}) :
  PD(\widehat{\fuk}(\lag, \lambda)) \longrightarrow
  PD(\widehat{\fuk}(\lag, \lambda)).
\end{equation}
The compatibility with respect to the comparison functors, as required
by diagrams~\eqref{sys-fun-diag-1} and~\eqref{sys-fun-diag-2}, follows
by standard arguments.

There is also a similar functor
$PD^c(\widehat{\phi}^{\fuk}) : PD^c(\widehat{\fuk}(\lag, \lambda))
\longrightarrow PD^c(\widehat{\fuk}(\lag, \lambda))$, defined
analogously to~\eqref{eq:phi-fun-3}, in which we just replace $PD$ by
$PD^c$ everywhere in~\eqref{eq:phi-fun-2}.


The definition of $\phi^{\fuk}_{p}$, $\widehat{\phi}^{\fuk}$ and their
persistence derived versions $PD(\phi^{\fuk}_{p})$,
$PD(\widehat{\phi}^{\fuk})$ in the monotone and the weakly exact cases
is very similar to the above. In these cases, the action on the
perturbation data $p \mapsto \phi_*(p)$ is precisely the same as in
the exact case described above. The action of $\phi$ on the grading in
the weakly exact case is the same as in the exact case, while in the
montone case it follows the receipe
from~\cite[Section~2.b]{Se:graded}. Finally, the action of $\phi$ on
the middle parameter $a_L \in \mathbb{R}$ of an object
$L = (\overline{L}, a_L, -)$ is the identity in both the monotone and
the weakly exact cases.


\subsection{Ungraded setting} \label{sb:fuk-ungraded} Sometimes it
will be more convenient to consider Fukaya categories without any
grading. This can be done in each of the settings described
in~\S\ref{sb:fuk-ex} -~\S\ref{sb:fuk-wex}, by simply omitting the
grading component from the objects $L \in \lag$ and viewing the hom's
as ungraded chain complexes. In the ungraded context, whenever we deal
with triangulated structures we just set the translation functor to be
$\id$. Note also that all the considerations from~\S\ref{sb:sys-fuk}
carry over to the ungraded setting in a straightforward way.
